\begin{otherlanguage}{french}
\chapter*{Résumé}
\addcontentsline{toc}{chapter}{Résumé}
\markboth{Résumé}{Résumé}

Chaque jour, des effets personnels --- téléphones, portefeuilles, papiers
d'identité, clés, sacs --- sont égarés dans l'espace public, et les moyens
permettant de les restituer à leurs propriétaires demeurent largement
informels. En Algérie, la récupération repose aujourd'hui sur des bureaux des
objets trouvés dépourvus de registre commun, ou sur des groupes de réseaux
sociaux où une annonce n'est qu'un texte non structuré qui disparaît du fil en
quelques heures. La conséquence est structurelle et non accidentelle: une
déclaration de perte et la déclaration de découverte correspondante peuvent
coexister, au même instant, sans jamais se rencontrer, faute d'un dispositif
chargé de les comparer.

Ce mémoire présente \projectname, une plateforme web qui automatise cette
comparaison. L'utilisateur déclare un objet perdu ou trouvé au moyen d'un
formulaire unique --- intitulé, description libre, catégorie, couleur, marque,
wilaya, lieu, date et jusqu'à cinq photographies --- et la plateforme interroge pour
lui l'ensemble opposé. L'appariement est \emph{multimodal}. Chaque déclaration
est encodée en un vecteur de phrase multilingue de dimension~384, de sorte
qu'une description rédigée en arabe, en français ou en anglais reste comparable
à une description rédigée dans l'une des deux autres langues~; chaque
photographie est encodée en un vecteur CLIP de dimension~512, de sorte que deux
personnes photographiant le même objet sous des angles différents produisent
des points voisins dans l'espace vectoriel. Les deux vecteurs sont stockés dans
PostgreSQL via l'extension \code{pgvector}, ce qui permet d'exprimer dans une
seule requête la recherche approximative des plus proches voisins et le
filtrage relationnel --- type de déclaration opposé, déclaration non encore
réglée, fenêtre temporelle plausible. Les candidats retenus sont ensuite évalués par une
fonction de fusion qui combine la similarité sémantique textuelle et un score
lexical plein texte, calibre la similarité visuelle brute, redistribue les
poids sur les seules modalités effectivement présentes, et ajoute de faibles
bonus bornés en cas de concordance de métadonnées indépendantes. Il en résulte
un indice de confiance, présenté à l'utilisateur sous forme de pourcentage
accompagné des raisons ayant motivé la suggestion.

Le système est réalisé sous la forme d'un monolithe modulaire: une API REST
versionnée développée avec FastAPI, un worker d'arrière-plan \code{arq} qui
concentre l'ensemble des inférences afin que la déclaration d'un objet
n'attende jamais un réseau de neurones, une base PostgreSQL~16 dotée de
\code{pgvector}, Redis comme file d'attente, et une interface Next.js
entièrement localisée en trois langues (anglais, français et arabe en écriture
de droite à gauche). Les coordonnées personnelles ne sont jamais divulguées par
un simple appariement: une procédure de réclamation distincte impose au
demandeur de répondre à des questions de vérification définies par le
déclarant, et l'échange d'identités n'intervient qu'après approbation. Le
prototype couvre le cycle complet --- inscription, déclaration, encodage,
appariement, notification, réclamation et restitution --- et chaque verdict
qu'un propriétaire rend sur une suggestion qu'il a consultée est conservé comme
exemple étiqueté en vue d'un futur modèle de classement calibré.

Autour de ce moteur, le prototype propose un appariement payant --- la
déclaration, la recherche et la réclamation restent gratuites, tandis que le
déblocage d'une suggestion au-delà de la première coûte un crédit, vendu par
lots de 200 à 1\,200 dinars via la passerelle de paiement Chargily --- ainsi
qu'une console d'administration dont chaque action est consignée dans un journal
d'audit au sein de la même transaction. Mesuré sur ses propres données,
l'encodeur multilingue a retrouvé le bon objet en première position pour les 72
requêtes d'une expérience interlangue contrôlée, contre 23 pour un modèle
uniquement anglophone~; la distribution mesurée des similarités visuelles, de
médiane 0{,}58 entre photographies sans rapport, a fixé le calibrage du canal
visuel~; et sept des huit suggestions jugées par leurs propriétaires étaient
correctes. L'ensemble de la plateforme fonctionne sur processeur avec environ
1{,}2\,Go de mémoire, ce qui permet d'exploiter un pilote public pour environ
4\,900 dinars par mois, charges du statut d'auto-entrepreneur comprises.

\vspace{0.8cm}
\noindent\textbf{Mots-clés:} objets perdus et trouvés, recherche sémantique,
plongements de phrases, CLIP, appariement multimodal, base de données
vectorielle, pgvector, recherche approximative des plus proches voisins,
FastAPI, Next.js, recherche d'information multilingue.
\end{otherlanguage}
